% ======================================================================
\chapter[Альфа-распад. Закон Гейгера--Нэттола. Формула Гамова]%
        {Альфа-распад. Закон Гейгера--Нэттола и его вывод (формула Гамова)}
\label{t:19}
% ======================================================================

\begin{qbox}
\textbf{Билет 19.} Альфа-распад, закон Гейгера--Нэттола и его вывод (формула Гамова).
\end{qbox}

\section{Содержание билета}

\textbf{Подтемы:}
\begin{enumerate}
 \item альфа-распад: схема, энергия распада, условие энергетической возможности, спектр альфа-частиц;
 \item экспериментальные закономерности: узкий интервал энергий и огромный разброс периодов; закон Гейгера--Нэттола;
 \item парадокс кулоновского барьера и его разрешение туннельным эффектом;
 \item теория Гамова: модель, вычисление проницаемости кулоновского барьера, формула для постоянной распада;
 \item вывод закона Гейгера--Нэттола; численные оценки и границы применимости.
\end{enumerate}

\textbf{Необходимые понятия:} квазиклассическая формула для коэффициента прохождения (\tref{t:07});
закон распада (\tref{t:18}); энергия связи и формула Вайцзеккера (\tref{t:17}); $e^2=1{,}44$~МэВ$\cdot$фм.

\section{Определения и физические принципы}

\begin{definition}[альфа-распад]
Самопроизвольное испускание ядром $\alpha$-частицы (ядра $\nucl{He}{4}$):
\begin{equation}\label{eq:19:scheme}
  \nuc{X}{A}{Z}\ \to\ \nuc{Y}{A-4}{Z-2}+\nuc{He}{4}{2}.
\end{equation}
\emph{Энергия распада} (энерговыделение)
\begin{equation}\label{eq:19:Q}
  Q_\alpha=\bigl[M_{\text{ат}}(A,Z)-M_{\text{ат}}(A-4,Z-2)-M_{\text{ат}}(\nucl{He}{4})\bigr]c^2=E_{\text{св}}(\alpha)+E_{\text{св}}(A-4,Z-2)-E_{\text{св}}(A,Z).
\end{equation}
Распад возможен при $Q_\alpha>0$.
\end{definition}

\begin{principle}{Закон Гейгера--Нэттола}
Для ядер одного радиоактивного ряда (в современной форме — для чётно-чётных ядер с одинаковым $Z$) логарифм периода
полураспада линейно зависит от $Q_\alpha^{-1/2}$:
\begin{equation}\label{eq:19:GN}
  \lg T_{1/2}=a\,\frac{Z}{\sqrt{Q_\alpha}}+b ,
\end{equation}
где $Z$ — заряд дочернего ядра, $a\approx1{,}6$--$1{,}7$ (при $Q$ в МэВ), $b$ слабо зависит от $Z$. Исторически (1911) закон
формулировался как линейная связь $\lg\lambda$ и логарифма пробега $\alpha$-частиц в воздухе.
\end{principle}

\emph{Физический смысл.} Небольшое изменение энергии $\alpha$-частицы меняет время жизни на много порядков: при
$Q_\alpha$ от 4 до 9~МэВ периоды полураспада меняются от $10^{10}$ лет до долей микросекунды. Такая зависимость характерна
для туннельного прохождения барьера: $\lambda\propto\exp(-\text{const}\cdot Z/\sqrt Q)$.

\emph{Область применимости теории Гамова.} Одночастичная модель: $\alpha$-частица считается «готовой» внутри ядра
(вероятность её формирования — множитель порядка $10^{-1}$--$10^{-2}$), ядро сферическое, орбитальный момент вылетающей
частицы $l=0$ (центробежный барьер не учитывается), квазиклассическое приближение ($D\ll1$). Теория предсказывает
$T_{1/2}$ с точностью до 1--2 порядков — достаточно, чтобы объяснить разброс на 24 порядка.

\section{Основная теория}

\subsection{Энергетика альфа-распада}

\paragraph{Условие распада.} $\alpha$-частица очень сильно связана ($E_{\text{св}}=28{,}3$~МэВ, т.\,е. 7,07~МэВ/нуклон). Из~\eqref{eq:19:Q}
$Q_\alpha\approx E_{\text{св}}(\alpha)-4\varepsilon_{\text{эфф}}$, где $\varepsilon_{\text{эфф}}=\Delta E_{\text{св}}/\Delta A$ — энергия связи нуклонов, «отрываемых» от
тяжёлого ядра. Поскольку для $A\gtrsim60$ удельная энергия связи убывает с ростом $A$ (главным образом из-за кулоновского
отталкивания), для достаточно тяжёлых ядер $4\varepsilon_{\text{эфф}}<28{,}3$~МэВ. Формула Вайцзеккера даёт $Q_\alpha>0$ при $A\gtrsim150$;
практически заметный $\alpha$-распад наблюдается у ядер тяжелее свинца (а также у некоторых редкоземельных ядер,
например $\nucl{Sm}{147}$).

\paragraph{Распределение энергии.} Распад двухчастичный, поэтому $\alpha$-частицы \emph{моноэнергетичны} (вывод~\ref{der:19:Talpha}):
\begin{equation}\label{eq:19:Ta}
  T_\alpha=Q_\alpha\frac{M_Y}{M_Y+M_\alpha}\approx Q_\alpha\frac{A-4}{A},
\end{equation}
остальное ($\approx2\,\%$) уносит ядро отдачи. Если распад идёт и на возбуждённые уровни дочернего ядра, в спектре появляются
несколько линий (тонкая структура $\alpha$-спектра), а последующие $\gamma$-кванты уносят разность энергий. Типичные $T_\alpha=4$--$9$~МэВ.

\subsection{Парадокс кулоновского барьера}

Вне ядра на $\alpha$-частицу действует кулоновское отталкивание дочернего ядра, внутри — сильное притяжение
(рис.~\ref{fig:19:pot}). Высота кулоновского барьера при радиусе ядра $R$:
\begin{equation}\label{eq:19:barrier}
  U_B=\frac{2Ze^2}{R}.
\end{equation}
Для урана ($Z=90$, $R\approx9$~фм): $U_B\approx2\cdot90\cdot1{,}44/9\approx29$~МэВ. Опыты Резерфорда по рассеянию $\alpha$-частиц с энергией
$8{,}8$~МэВ на уране не обнаружили отклонений от кулоновского закона — барьер действительно выше 8,8~МэВ. Но сам $\nucl{U}{238}$
испускает $\alpha$-частицы с энергией лишь $4{,}2$~МэВ. Классически частица с энергией $T<U_B$ не может преодолеть барьер ни
изнутри, ни снаружи. Объяснение дали Г.~Гамов и независимо Р.~Гёрни и Э.~Кондон (1928): $\alpha$-частица \emph{туннелирует}
сквозь барьер. Это было первое применение квантовой механики к ядру.

\begin{figure}[htbp]
\centering
\begin{tikzpicture}
\begin{axis}[width=0.8\linewidth, height=6.0cm, axis lines=middle, xlabel={$r$}, ylabel={$U(r)$},
   xmin=0, xmax=58, ymin=-42, ymax=34, xtick={9,46}, xticklabels={$R$,$b$}, ytick={5.6,28.8,-35},
   yticklabels={$Q_\alpha$,$U_B=\frac{2Ze^2}{R}$,$-U_0$}, samples=200, clip=false,
   xlabel style={at={(ticklabel* cs:1)}, anchor=west}, ylabel style={at={(ticklabel* cs:1)}, anchor=south},
   yticklabel style={font=\small}]
  \addplot[very thick] coordinates {(0,-35) (9,-35) (9,28.8)};
  \addplot[very thick, domain=9:57, name path=C] {259.2/x};
  \addplot[gray, dashed, domain=0:57, name path=E] {5.6};
  \addplot[accent!30] fill between[of=C and E, soft clip={domain=9:46.3}];
  \node[font=\small, accent, anchor=west] at (axis cs:15,8.2) {туннелирование};
  \draw[->, keycol, thick] (axis cs:47,5.6) -- (axis cs:56,5.6);
  \node[font=\small, keycol, anchor=south] at (axis cs:52,6) {$\alpha$};
  \node[font=\small, anchor=west] at (axis cs:28,24) {$U=\dfrac{2Ze^2}{r}$};
  \node[font=\small, align=center] at (axis cs:4.5,-28) {ядерная\\яма};
\end{axis}
\end{tikzpicture}
\caption{Потенциальная энергия $\alpha$-частицы в поле дочернего ядра (модель Гамова): прямоугольная яма глубины $U_0$
при $r<R$ и кулоновский потенциал $2Ze^2/r$ при $r>R$. Под барьером закрашена область $R<r<b$, $b=2Ze^2/Q_\alpha$ —
внешняя точка поворота. Масштаб условный: для $\nucl{U}{238}$ $Q_\alpha\approx4{,}3$~МэВ, $U_B\approx29$~МэВ, $b\approx7R$.}
\label{fig:19:pot}
\end{figure}

\subsection{Теория Гамова}

\paragraph{Модель.} (1) $\alpha$-частица предсуществует в ядре и движется в потенциальной яме радиуса $R$, ударяясь о стенку с
частотой $\nu_0\approx v_{\text{вн}}/2R$; (2) при каждом ударе она проходит барьер с вероятностью $D$; (3) постоянная распада
\begin{equation}\label{eq:19:lambda}
  \lambda=\nu_0D .
\end{equation}
Оценка: скорость внутри ядра $v_{\text{вн}}\sim0{,}1c$ (кинетическая энергия $\sim Q+U_0\sim$ десятки МэВ), $R\sim10$~фм:
$\nu_0\sim3\cdot10^{22}\ \text{фм/с}\cdot0{,}1/(20\ \text{фм})\sim10^{21}$~с$^{-1}$.

\paragraph{Проницаемость барьера.} По квазиклассической формуле (\tref{t:07})
\[
  D\approx\exp\Bigl[-\frac2\hbar\int_R^b\sqrt{2m\Bigl(\frac{2Ze^2}{r}-Q\Bigr)}\,\dd r\Bigr],\qquad b=\frac{2Ze^2}{Q},
\]
где $m$ — приведённая масса $\alpha$-частицы и ядра ($\approx m_\alpha$), $Q\approx T_\alpha$. Вычисление интеграла
(вывод~\ref{der:19:gamow}) даёт \emph{формулу Гамова}
\begin{equation}\label{eq:19:Gamow}
  D\approx\exp\Bigl[-\frac{4\pi Ze^2}{\hbar v}+\frac{8}{\hbar}\sqrt{mZe^2R}\Bigr],\qquad v=\sqrt{\frac{2Q}{m}},
\end{equation}
где $v$ — скорость $\alpha$-частицы на бесконечности. Главный член $2\pi\eta=4\pi Ze^2/\hbar v$ («фактор Гамова») зависит только от
заряда и скорости; второй — поправка на конечный радиус ядра (барьер «срезан» при $r=R$).

\paragraph{Постоянная распада:}
\begin{equation}\label{eq:19:lnlambda}
  \ln\lambda=\ln\nu_0+\frac8\hbar\sqrt{mZe^2R}-\frac{4\pi Ze^2}{\hbar}\sqrt{\frac{m}{2Q}}\,.
\end{equation}

\subsection{Вывод закона Гейгера--Нэттола}

В~\eqref{eq:19:lnlambda} последний член зависит от энергии как $Z/\sqrt Q$, а первые два слабо (логарифмически или через $\sqrt{ZR}$)
зависят от $Z$ и $A$. Для ядер с близкими $Z$ и $A$ (один ряд, или фиксированное $Z$) они практически постоянны. Переходя к
$T_{1/2}=\ln2/\lambda$ и десятичным логарифмам:
\begin{equation}\label{eq:19:GNderived}
  \lg T_{1/2}=\frac{4\pi e^2}{\hbar\ln10}\sqrt{\frac m2}\cdot\frac{Z}{\sqrt Q}+b(Z,R)
  \approx1{,}7\,\frac{Z}{\sqrt{Q\,[\text{МэВ}]}}+b ,
\end{equation}
что и есть закон Гейгера--Нэттола~\eqref{eq:19:GN}, причём коэффициент $a$ выражен через фундаментальные константы
(вывод~\ref{der:19:coef}). Обратно: из сравнения экспериментальных $T_{1/2}$ с~\eqref{eq:19:lnlambda} получают радиусы ядер
$R\approx1{,}4$--$1{,}5\,A^{1/3}$~фм — один из первых методов их определения.

\begin{figure}[htbp]
\centering
\begin{tikzpicture}
\begin{axis}[width=0.7\linewidth, height=5.6cm, xlabel={$Q_\alpha$, МэВ}, ylabel={$\lg T_{1/2}$ [с]}, xmin=4, xmax=9.2,
   ymin=-9, ymax=21, samples=100, grid=major, grid style={gray!20},
   legend style={font=\small, draw=none}, legend pos=north east,
   declare function={lgT(\Q,\Z)=(0.0917*\Z*43.17/sqrt(\Q)-2.97*sqrt(\Z*1.45*(2.5*\Z)^(1/3)))/2.3026-21.16;}]
  \addplot[very thick, accent, domain=4:9.2] {lgT(x,82)};
  \addlegendentry{$Z=82$}
  \addplot[very thick, keycol, dashed, domain=4:9.2] {lgT(x,86)};
  \addlegendentry{$Z=86$}
  \addplot[very thick, dercol, dotted, domain=4:9.2] {lgT(x,90)};
  \addlegendentry{$Z=90$}
\end{axis}
\end{tikzpicture}
\caption{Оценка периода полураспада по формуле Гамова~\eqref{eq:19:lnlambda} для разных зарядов дочернего ядра
($R=1{,}45A_Y^{1/3}$~фм, $A_Y\approx2{,}5Z$, $\nu_0=10^{21}$~с$^{-1}$; точность оценки — 1--2 порядка). Изменение $Q_\alpha$ вдвое
меняет $T_{1/2}$ более чем на 20 порядков.}
\label{fig:19:GN}
\end{figure}

\paragraph{Факторы, не учтённые простой теорией.} Вероятность образования $\alpha$-кластера в ядре (уменьшает $\lambda$ в 10--100 раз);
центробежный барьер $\hbar^2l(l+1)/2mr^2$ при вылете с $l\neq0$ (замедляет распад); деформация ядра (несферический барьер);
для нечётных ядер распады часто «затруднены» (структурные эффекты). Учёт этих факторов улучшает согласие до множителя
порядка единицы.

\section{Математические выводы}

\begin{derivation}[label=der:19:Talpha]{кинетическая энергия альфа-частицы}
Материнское ядро покоится. Законы сохранения (нерелятивистски, $Q\ll Mc^2$):
\[
  p_\alpha=p_Y\equiv p,\qquad T_\alpha+T_Y=\frac{p^2}{2M_\alpha}+\frac{p^2}{2M_Y}=Q .
\]
Отсюда $\dfrac{T_\alpha}{T_Y}=\dfrac{M_Y}{M_\alpha}$ и
\[
  T_\alpha=Q\frac{M_Y}{M_\alpha+M_Y}\approx Q\frac{A-4}{A},\qquad T_Y=Q\frac{4}{A}.
\]
Для $\nucl{U}{238}$ ($Q=4{,}27$~МэВ): $T_\alpha=4{,}20$~МэВ, $T_Y=0{,}07$~МэВ. Энергия отдачи $\sim70$~кэВ — много больше энергии
химических связей, поэтому дочерние атомы выбиваются из молекул (используется в радиохимии).
\end{derivation}

\begin{derivation}[label=der:19:gamow]{формула Гамова}
\emph{Допущения:} квазиклассика (показатель $\gg1$); $l=0$; барьер — чисто кулоновский при $r>R$; ядро отдачи бесконечно тяжёлое
($m$ — приведённая масса).

\emph{Шаг 1 (интеграл).} Обозначим $b=2Ze^2/Q$ — внешнюю точку поворота. Тогда $\frac{2Ze^2}{r}-Q=Q\bigl(\frac br-1\bigr)$ и
\[
  I\equiv\int_R^b\sqrt{2m\Bigl(\frac{2Ze^2}{r}-Q\Bigr)}\dd r=\sqrt{2mQ}\int_R^b\sqrt{\frac br-1}\,\dd r .
\]
\emph{Шаг 2 (замена $r=b\cos^2u$).} $\dd r=-2b\cos u\sin u\,\dd u$, $\sqrt{b/r-1}=\tg u$. Пределы: $r=b\to u=0$; $r=R\to u_0=\arccos\sqrt{R/b}$:
\[
  \int_R^b\sqrt{\frac br-1}\,\dd r=\int_0^{u_0}\tg u\cdot2b\cos u\sin u\,\dd u=2b\int_0^{u_0}\sin^2u\,\dd u
  =b\bigl[u-\sin u\cos u\bigr]_0^{u_0}.
\]
С $\cos u_0=\sqrt{x}$, $\sin u_0=\sqrt{1-x}$, $x=R/b$:
\[
  I=\sqrt{2mQ}\,b\Bigl[\arccos\sqrt{x}-\sqrt{x(1-x)}\Bigr].
\]
\emph{Шаг 3 (разложение при $x=R/b\ll1$).} Для урана $x\approx0{,}15$. $\arccos\sqrt x\approx\frac\pi2-\sqrt x$, $\sqrt{x(1-x)}\approx\sqrt x$:
\[
  I\approx\sqrt{2mQ}\,b\Bigl(\frac\pi2-2\sqrt{\frac Rb}\Bigr).
\]
\emph{Шаг 4 (подстановка $b$).}
\[
  \frac2\hbar\sqrt{2mQ}\,b\,\frac\pi2=\frac{\pi}{\hbar}\sqrt{2mQ}\frac{2Ze^2}{Q}=\frac{2\pi Ze^2}{\hbar}\sqrt{\frac{2m}{Q}}=\frac{4\pi Ze^2}{\hbar v},
  \quad v=\sqrt{\frac{2Q}{m}};
\]
\[
  \frac2\hbar\sqrt{2mQ}\,b\cdot2\sqrt{\frac Rb}=\frac4\hbar\sqrt{2mQ\,Rb}=\frac4\hbar\sqrt{2mQR\frac{2Ze^2}{Q}}=\frac8\hbar\sqrt{mZe^2R}.
\]
\emph{Итог:} $D=\exp(-2I/\hbar)\approx\exp\Bigl[-\dfrac{4\pi Ze^2}{\hbar v}+\dfrac8\hbar\sqrt{mZe^2R}\Bigr]$.

\emph{Размерности:} $[Ze^2/\hbar v]=\dfrac{\text{эрг}\cdot\text{см}}{\text{эрг}\cdot\text{с}\cdot\text{см}/\text{с}}=1$; $[\sqrt{me^2R}/\hbar]=\dfrac{\sqrt{\text{г}\cdot\text{эрг}\,\text{см}\cdot\text{см}}}{\text{эрг}\,\text{с}}=\dfrac{\text{г}\,\text{см}^2/\text{с}}{\text{г}\,\text{см}^2/\text{с}}=1$.
\end{derivation}

\begin{derivation}[label=der:19:coef]{коэффициент закона Гейгера--Нэттола}
В~\eqref{eq:19:lnlambda} член, зависящий от энергии: $-\dfrac{4\pi Ze^2}{\hbar c}\sqrt{\dfrac{mc^2}{2Q}}=-4\pi\alpha Z\sqrt{\dfrac{mc^2}{2Q}}$. С $mc^2\approx3727$~МэВ:
$\sqrt{mc^2/2}=43{,}2\ \text{МэВ}^{1/2}$, $4\pi\alpha\cdot43{,}2=0{,}0917\cdot43{,}2=3{,}96\ \text{МэВ}^{1/2}$. Тогда
\[
  \ln T_{1/2}=-\ln\lambda+\ln\ln2=3{,}96\frac{Z}{\sqrt{Q[\text{МэВ}]}}+\dots,\qquad
  \lg T_{1/2}=\frac{3{,}96}{2{,}303}\frac{Z}{\sqrt Q}+\dots=1{,}72\frac{Z}{\sqrt Q}+b .
\]
Эмпирический коэффициент $\approx1{,}6$--$1{,}7$ — хорошее согласие (с учётом приведённой массы $m\approx0{,}98m_\alpha$ теоретический
коэффициент $\approx1{,}70$).
\end{derivation}

\begin{derivation}[label=der:19:num]{численная оценка для $\nucl{U}{238}$ и сравнение с $\nucl{Po}{212}$}
\emph{$\nucl{U}{238}\to\nucl{Th}{234}$:} $Z=90$, $Q=4{,}27$~МэВ, $R\approx1{,}45\cdot234^{1/3}\approx8{,}9$~фм.
Первый член: $3{,}96\cdot90/\sqrt{4{,}27}\approx172$. Второй: $\dfrac{8}{\hbar c}\sqrt{mc^2\,Ze^2R}=\dfrac{8\sqrt{3727\cdot90\cdot1{,}44\cdot8{,}9}}{197{,}3}\approx84$.
$\ln D\approx-172+84=-88$, $D\approx4\cdot10^{-39}$. При $\nu_0\sim10^{21}$~с$^{-1}$: $\lambda\sim4\cdot10^{-18}$~с$^{-1}$,
$T_{1/2}\sim2\cdot10^{17}$~с $\sim5\cdot10^9$~лет (опыт: $4{,}47\cdot10^9$~лет). Совпадение по порядку отчасти удачно: изменение $R$ на 10\,\%
меняет второй член на $\approx4$, т.\,е. $T_{1/2}$ — в $\ee^4\approx50$ раз.

\emph{Сравнение:} $\nucl{Po}{212}\to\nucl{Pb}{208}$ ($Z=82$, $Q=8{,}95$~МэВ) имеет $Z/\sqrt Q=27{,}4$ против $43{,}6$ для $\nucl{Th}{232}$ ($Z=88$, $Q=4{,}08$~МэВ).
Разность $16{,}2$ даёт по главному члену $1{,}72\cdot16{,}2\approx28$ порядков; с учётом поправки на радиус — около 26. Опыт:
$T_{1/2}(\nucl{Th}{232})=1{,}4\cdot10^{10}$~лет $=4{,}4\cdot10^{17}$~с, $T_{1/2}(\nucl{Po}{212})=3\cdot10^{-7}$~с — различие в 24 порядка при различии энергий
лишь вдвое.
\end{derivation}

\section{Итоговые формулы}

\keyformula{Q_\alpha=\bigl[M(A,Z)-M(A-4,Z-2)-M(\nucl{He}{4})\bigr]c^2,\qquad T_\alpha\approx Q_\alpha\frac{A-4}{A}}
 {$M$ — массы атомов; $T_\alpha$ — кинетическая энергия $\alpha$-частицы.}
 {распад из основного состояния в основное.}
 {условие распада ($Q>0$), энергия линий $\alpha$-спектра.}

\keyformula{D\approx\exp\Bigl[-\frac{4\pi Ze^2}{\hbar v}+\frac8\hbar\sqrt{mZe^2R}\Bigr],\qquad \lambda=\nu_0D}
 {$Z$ — заряд дочернего ядра, $v=\sqrt{2Q/m}$, $m\approx m_\alpha$, $R$ — радиус ядра, $\nu_0\sim10^{21}$~с$^{-1}$.}
 {квазиклассика, $R\ll b=2Ze^2/Q$, $l=0$, предсуществование $\alpha$-частицы.}
 {постоянная $\alpha$-распада; радиусы ядер по периодам полураспада.}

\keyformula{\lg T_{1/2}=a\frac{Z}{\sqrt{Q_\alpha}}+b,\qquad a=\frac{4\pi\alpha}{\ln10}\sqrt{\frac{mc^2}{2}}\approx1{,}7\ \text{МэВ}^{1/2}}
 {$Q_\alpha$ в МэВ; $b$ слабо зависит от $Z$ и $R$.}
 {ядра с близкими $Z$, $A$; чётно-чётные ядра, $l=0$.}
 {закон Гейгера--Нэттола; экспоненциальная чувствительность к энергии.}

\section{Физическая интерпретация}

\textbf{Почему периоды так различны.} Показатель экспоненты $\sim\pi Ze^2\cdot2/\hbar v\sim100$; изменение энергии на 10\,\% меняет его на 5\,\%,
т.\,е. на несколько единиц — и $T_{1/2}$ на несколько порядков. Поэтому $\alpha$-активные ядра с $Q<4$~МэВ практически стабильны
($T\gg$ возраста Вселенной), а с $Q>9$~МэВ живут микросекунды. Узкий «наблюдаемый» интервал энергий $4$--$9$~МэВ — следствие
этой зависимости.

\textbf{Почему испускаются $\alpha$-частицы, а не протоны или $\nucl{C}{12}$.} Испускание отдельного нуклона обычно энергетически
запрещено ($Q<0$), а $\alpha$-частица очень сильно связана и даёт $Q>0$. Более тяжёлые кластеры ($\nucl{C}{14}$, $\nucl{Ne}{24}$) иногда
энергетически разрешены, но их барьер ($\propto Z_1Z_2$) и масса делают проницаемость ничтожной: кластерный распад наблюдается с
вероятностью $\sim10^{-9}$--$10^{-16}$ относительно $\alpha$-распада.

\textbf{Обращённый процесс.} Тот же фактор Гамова $\exp(-2\pi\eta)$ определяет вероятность сближения заряженных ядер при
термоядерных реакциях в звёздах: синтез при $kT\sim1$~кэВ $\ll$ кулоновского барьера ($\sim1$~МэВ) идёт только благодаря туннелированию.

\section{Что необходимо знать на экзамене}

\begin{examcheck}
\checkhead{Определения}
\begin{itemize}
 \item схема $\alpha$-распада, $Q_\alpha$, условие $Q_\alpha>0$; моноэнергетичность и тонкая структура спектра;
 \item закон Гейгера--Нэттола (формулировка, смысл коэффициентов).
\end{itemize}
\checkhead{Формулы}
\begin{itemize}
 \item $T_\alpha=Q(A-4)/A$; $U_B=2Ze^2/R$; $\lambda=\nu_0D$;
 \item $D\approx\exp[-4\pi Ze^2/\hbar v+(8/\hbar)\sqrt{mZe^2R}]$; $\lg T=aZ/\sqrt Q+b$, $a\approx1{,}7$.
\end{itemize}
\checkhead{Выводы}
\begin{itemize}
 \item парадокс барьера (рассеяние 8,8~МэВ vs испускание 4,2~МэВ) и его разрешение;
 \item вычисление интеграла Гамова (замена $r=b\cos^2u$, разложение при $R\ll b$);
 \item переход к закону Гейгера--Нэттола и оценка коэффициента $a$.
\end{itemize}
\checkhead{Интерпретации}
\begin{itemize}
 \item модель предсуществующей частицы и частоты ударов; неучтённые факторы (образование кластера, $l\neq0$, деформация);
 \item экспоненциальная чувствительность $T_{1/2}$ к $Q$ и $R$; связь с термоядерным синтезом.
\end{itemize}
\end{examcheck}
